\IfFileExists{stacks-project.cls}{%
\documentclass{stacks-project}
}{%
\documentclass{amsart}
}

% For dealing with references we use the comment environment
\usepackage{verbatim}
\newenvironment{reference}{\comment}{\endcomment}
%\newenvironment{reference}{}{}
\newenvironment{slogan}{\comment}{\endcomment}
\newenvironment{history}{\comment}{\endcomment}

% For commutative diagrams we use Xy-pic
\usepackage[all]{xy}

% We use 2cell for 2-commutative diagrams.
\xyoption{2cell}
\UseAllTwocells

% We use multicol for the list of chapters between chapters
\usepackage{multicol}

% This is generally recommended for better output
\usepackage{lmodern}
\usepackage[T1]{fontenc}

% For cross-file-references
\usepackage{xr-hyper}

% Package for hypertext links:
\usepackage{hyperref}

% For any local file, say "hello.tex" you want to link to please
% use \externaldocument[hello-]{hello}
\externaldocument[introduction-]{introduction}
\externaldocument[conventions-]{conventions}
\externaldocument[sets-]{sets}
\externaldocument[categories-]{categories}
\externaldocument[topology-]{topology}
\externaldocument[sheaves-]{sheaves}
\externaldocument[sites-]{sites}
\externaldocument[stacks-]{stacks}
\externaldocument[fields-]{fields}
\externaldocument[algebra-]{algebra}
\externaldocument[brauer-]{brauer}
\externaldocument[homology-]{homology}
\externaldocument[derived-]{derived}
\externaldocument[simplicial-]{simplicial}
\externaldocument[more-algebra-]{more-algebra}
\externaldocument[smoothing-]{smoothing}
\externaldocument[modules-]{modules}
\externaldocument[sites-modules-]{sites-modules}
\externaldocument[injectives-]{injectives}
\externaldocument[cohomology-]{cohomology}
\externaldocument[sites-cohomology-]{sites-cohomology}
\externaldocument[dga-]{dga}
\externaldocument[dpa-]{dpa}
\externaldocument[sdga-]{sdga}
\externaldocument[hypercovering-]{hypercovering}
\externaldocument[schemes-]{schemes}
\externaldocument[constructions-]{constructions}
\externaldocument[properties-]{properties}
\externaldocument[morphisms-]{morphisms}
\externaldocument[coherent-]{coherent}
\externaldocument[divisors-]{divisors}
\externaldocument[limits-]{limits}
\externaldocument[varieties-]{varieties}
\externaldocument[topologies-]{topologies}
\externaldocument[descent-]{descent}
\externaldocument[perfect-]{perfect}
\externaldocument[more-morphisms-]{more-morphisms}
\externaldocument[flat-]{flat}
\externaldocument[groupoids-]{groupoids}
\externaldocument[more-groupoids-]{more-groupoids}
\externaldocument[etale-]{etale}
\externaldocument[chow-]{chow}
\externaldocument[intersection-]{intersection}
\externaldocument[pic-]{pic}
\externaldocument[weil-]{weil}
\externaldocument[adequate-]{adequate}
\externaldocument[dualizing-]{dualizing}
\externaldocument[duality-]{duality}
\externaldocument[discriminant-]{discriminant}
\externaldocument[derham-]{derham}
\externaldocument[local-cohomology-]{local-cohomology}
\externaldocument[algebraization-]{algebraization}
\externaldocument[curves-]{curves}
\externaldocument[resolve-]{resolve}
\externaldocument[models-]{models}
\externaldocument[functors-]{functors}
\externaldocument[equiv-]{equiv}
\externaldocument[pione-]{pione}
\externaldocument[etale-cohomology-]{etale-cohomology}
\externaldocument[proetale-]{proetale}
\externaldocument[relative-cycles-]{relative-cycles}
\externaldocument[more-etale-]{more-etale}
\externaldocument[trace-]{trace}
\externaldocument[crystalline-]{crystalline}
\externaldocument[spaces-]{spaces}
\externaldocument[spaces-properties-]{spaces-properties}
\externaldocument[spaces-morphisms-]{spaces-morphisms}
\externaldocument[decent-spaces-]{decent-spaces}
\externaldocument[spaces-cohomology-]{spaces-cohomology}
\externaldocument[spaces-limits-]{spaces-limits}
\externaldocument[spaces-divisors-]{spaces-divisors}
\externaldocument[spaces-over-fields-]{spaces-over-fields}
\externaldocument[spaces-topologies-]{spaces-topologies}
\externaldocument[spaces-descent-]{spaces-descent}
\externaldocument[spaces-perfect-]{spaces-perfect}
\externaldocument[spaces-more-morphisms-]{spaces-more-morphisms}
\externaldocument[spaces-flat-]{spaces-flat}
\externaldocument[spaces-groupoids-]{spaces-groupoids}
\externaldocument[spaces-more-groupoids-]{spaces-more-groupoids}
\externaldocument[bootstrap-]{bootstrap}
\externaldocument[spaces-pushouts-]{spaces-pushouts}
\externaldocument[spaces-chow-]{spaces-chow}
\externaldocument[groupoids-quotients-]{groupoids-quotients}
\externaldocument[spaces-more-cohomology-]{spaces-more-cohomology}
\externaldocument[spaces-simplicial-]{spaces-simplicial}
\externaldocument[spaces-duality-]{spaces-duality}
\externaldocument[formal-spaces-]{formal-spaces}
\externaldocument[restricted-]{restricted}
\externaldocument[spaces-resolve-]{spaces-resolve}
\externaldocument[formal-defos-]{formal-defos}
\externaldocument[defos-]{defos}
\externaldocument[cotangent-]{cotangent}
\externaldocument[examples-defos-]{examples-defos}
\externaldocument[algebraic-]{algebraic}
\externaldocument[examples-stacks-]{examples-stacks}
\externaldocument[stacks-sheaves-]{stacks-sheaves}
\externaldocument[criteria-]{criteria}
\externaldocument[artin-]{artin}
\externaldocument[quot-]{quot}
\externaldocument[stacks-properties-]{stacks-properties}
\externaldocument[stacks-morphisms-]{stacks-morphisms}
\externaldocument[stacks-limits-]{stacks-limits}
\externaldocument[stacks-cohomology-]{stacks-cohomology}
\externaldocument[stacks-perfect-]{stacks-perfect}
\externaldocument[stacks-introduction-]{stacks-introduction}
\externaldocument[stacks-more-morphisms-]{stacks-more-morphisms}
\externaldocument[stacks-geometry-]{stacks-geometry}
\externaldocument[moduli-]{moduli}
\externaldocument[moduli-curves-]{moduli-curves}
\externaldocument[examples-]{examples}
\externaldocument[exercises-]{exercises}
\externaldocument[guide-]{guide}
\externaldocument[desirables-]{desirables}
\externaldocument[coding-]{coding}
\externaldocument[obsolete-]{obsolete}
\externaldocument[fdl-]{fdl}
\externaldocument[index-]{index}

% Theorem environments.
%
\theoremstyle{plain}
\newtheorem{theorem}[subsection]{Theorem}
\newtheorem{proposition}[subsection]{Proposition}
\newtheorem{lemma}[subsection]{Lemma}

\theoremstyle{definition}
\newtheorem{definition}[subsection]{Definition}
\newtheorem{example}[subsection]{Example}
\newtheorem{exercise}[subsection]{Exercise}
\newtheorem{situation}[subsection]{Situation}

\theoremstyle{remark}
\newtheorem{remark}[subsection]{Remark}
\newtheorem{remarks}[subsection]{Remarks}

\numberwithin{equation}{subsection}

% Macros
%
\def\lim{\mathop{\mathrm{lim}}\nolimits}
\def\colim{\mathop{\mathrm{colim}}\nolimits}
\def\Spec{\mathop{\mathrm{Spec}}}
\def\Hom{\mathop{\mathrm{Hom}}\nolimits}
\def\Ext{\mathop{\mathrm{Ext}}\nolimits}
\def\SheafHom{\mathop{\mathcal{H}\!\mathit{om}}\nolimits}
\def\SheafExt{\mathop{\mathcal{E}\!\mathit{xt}}\nolimits}
\def\Sch{\mathit{Sch}}
\def\Mor{\mathop{\mathrm{Mor}}\nolimits}
\def\Ob{\mathop{\mathrm{Ob}}\nolimits}
\def\Sh{\mathop{\mathit{Sh}}\nolimits}
\def\NL{\mathop{N\!L}\nolimits}
\def\CH{\mathop{\mathrm{CH}}\nolimits}
\def\proetale{{pro\text{-}\acute{e}tale}}
\def\etale{{\acute{e}tale}}
\def\QCoh{\mathit{QCoh}}
\def\Ker{\mathop{\mathrm{Ker}}}
\def\Im{\mathop{\mathrm{Im}}}
\def\Coker{\mathop{\mathrm{Coker}}}
\def\Coim{\mathop{\mathrm{Coim}}}

% Boxtimes
%
\DeclareMathSymbol{\boxtimes}{\mathbin}{AMSa}{"02}

%
% Macros for moduli stacks/spaces
%
\def\QCohstack{\mathcal{QC}\!\mathit{oh}}
\def\Cohstack{\mathcal{C}\!\mathit{oh}}
\def\Spacesstack{\mathcal{S}\!\mathit{paces}}
\def\Quotfunctor{\mathrm{Quot}}
\def\Hilbfunctor{\mathrm{Hilb}}
\def\Curvesstack{\mathcal{C}\!\mathit{urves}}
\def\Polarizedstack{\mathcal{P}\!\mathit{olarized}}
\def\Complexesstack{\mathcal{C}\!\mathit{omplexes}}
% \Pic is the operator that assigns to X its picard group, usage \Pic(X)
% \Picardstack_{X/B} denotes the Picard stack of X over B
% \Picardfunctor_{X/B} denotes the Picard functor of X over B
\def\Pic{\mathop{\mathrm{Pic}}\nolimits}
\def\Picardstack{\mathcal{P}\!\mathit{ic}}
\def\Picardfunctor{\mathrm{Pic}}
\def\Deformationcategory{\mathcal{D}\!\mathit{ef}}
\usepackage{cleveref}
\crefname{originallemma}{original lemma}{original lemmas}
\Crefname{originallemma}{Original lemma}{Original lemmas}
\makeatletter
\newlabel{coherent-lemma-quasi-coherent-affine-cohomology-zero}{{01XB}{source}{}{sourcepointer-1}{}}
\newlabel{coherent-lemma-quasi-coherent-affine-cohomology-zero@cref}{{[originallemma][0][]01XB}{[1][0][]source}{}{}{}}
\makeatother
\crefname{originalremark}{original remark}{original remarks}
\Crefname{originalremark}{Original remark}{Original remarks}
\makeatletter
\newlabel{hypercovering-remark-P-covering}{{0DER}{source}{}{sourcepointer-2}{}}
\newlabel{hypercovering-remark-P-covering@cref}{{[originalremark][0][]0DER}{[1][0][]source}{}{}{}}
\makeatother
\crefname{originalreference}{original reference}{original references}
\Crefname{originalreference}{Original reference}{Original references}
\makeatletter
\newlabel{lemma-Serre-subcat-cartesian-modules}{{0D7J}{source}{}{sourcepointer-3}{}}
\newlabel{lemma-Serre-subcat-cartesian-modules@cref}{{[originalreference][0][]0D7J}{[1][0][]source}{}{}{}}
\makeatother
\crefname{originalreference}{original reference}{original references}
\Crefname{originalreference}{Original reference}{Original references}
\makeatletter
\newlabel{lemma-cohomological-descent-qcoh-for-fppf-hypercovering}{{0DHE}{source}{}{sourcepointer-4}{}}
\newlabel{lemma-cohomological-descent-qcoh-for-fppf-hypercovering@cref}{{[originalreference][0][]0DHE}{[1][0][]source}{}{}{}}
\makeatother
\crefname{originalreference}{original reference}{original references}
\Crefname{originalreference}{Original reference}{Original references}
\makeatletter
\newlabel{lemma-descent-qcoh-for-fppf-hypercovering}{{0DHD}{source}{}{sourcepointer-5}{}}
\newlabel{lemma-descent-qcoh-for-fppf-hypercovering@cref}{{[originalreference][0][]0DHD}{[1][0][]source}{}{}{}}
\makeatother
\crefname{originalreference}{original reference}{original references}
\Crefname{originalreference}{Original reference}{Original references}
\makeatletter
\newlabel{lemma-flat-augmentation-modules}{{0DH3}{source}{}{sourcepointer-6}{}}
\newlabel{lemma-flat-augmentation-modules@cref}{{[originalreference][0][]0DH3}{[1][0][]source}{}{}{}}
\makeatother
\crefname{originalreference}{original reference}{original references}
\Crefname{originalreference}{Original reference}{Original references}
\makeatletter
\newlabel{lemma-quasi-coherent-sheaf}{{0D7M}{source}{}{sourcepointer-7}{}}
\newlabel{lemma-quasi-coherent-sheaf@cref}{{[originalreference][0][]0D7M}{[1][0][]source}{}{}{}}
\makeatother
\crefname{originalreference}{original reference}{original references}
\Crefname{originalreference}{Original reference}{Original references}
\makeatletter
\newlabel{lemma-sanity-check-modules}{{0DH2}{source}{}{sourcepointer-8}{}}
\newlabel{lemma-sanity-check-modules@cref}{{[originalreference][0][]0DH2}{[1][0][]source}{}{}{}}
\makeatother
\crefname{originalsection}{original section}{original sections}
\Crefname{originalsection}{Original section}{Original sections}
\makeatletter
\newlabel{section-simplicial-algebraic-spaces}{{0DE7}{source}{}{sourcepointer-9}{}}
\newlabel{section-simplicial-algebraic-spaces@cref}{{[originalsection][0][]0DE7}{[1][0][]source}{}{}{}}
\makeatother
\crefname{originallemma}{original lemma}{original lemmas}
\Crefname{originallemma}{Original lemma}{Original lemmas}
\makeatletter
\newlabel{sites-cohomology-lemma-equivalence-unbounded-one}{{0D7V}{source}{}{sourcepointer-10}{}}
\newlabel{sites-cohomology-lemma-equivalence-unbounded-one@cref}{{[originallemma][0][]0D7V}{[1][0][]source}{}{}{}}
\makeatother
\crefname{originalreference}{original reference}{original references}
\Crefname{originalreference}{Original reference}{Original references}
\makeatletter
\newlabel{sites-cohomology-situation-olsson-laszlo}{{0D6R}{source}{}{sourcepointer-11}{}}
\newlabel{sites-cohomology-situation-olsson-laszlo@cref}{{[originalreference][0][]0D6R}{[1][0][]source}{}{}{}}
\makeatother
\crefname{originalsection}{original section}{original sections}
\Crefname{originalsection}{Original section}{Original sections}
\makeatletter
\newlabel{spaces-cohomology-section-higher-direct-image}{{071Y}{source}{}{sourcepointer-12}{}}
\newlabel{spaces-cohomology-section-higher-direct-image@cref}{{[originalsection][0][]071Y}{[1][0][]source}{}{}{}}
\makeatother
\crefname{originallemma}{original lemma}{original lemmas}
\Crefname{originallemma}{Original lemma}{Original lemmas}
\makeatletter
\newlabel{spaces-properties-lemma-properties-quasi-coherent}{{03M1}{source}{}{sourcepointer-13}{}}
\newlabel{spaces-properties-lemma-properties-quasi-coherent@cref}{{[originallemma][0][]03M1}{[1][0][]source}{}{}{}}
\makeatother
\title{fppf hypercoverings: QCoh is weak Serre and a\textasciicircum{}*: D\_QCoh(O\_X) -> D\_QCoh(O\_U) is an equivalence with quasi-inverse Ra\_*}
\author{The Stacks Project Authors}
\date{stacks-project:spaces-simplicial.tex:local-snapshot-20261012; GNU FDL 1.2 or later}
\begin{document}
\maketitle
\noindent\textit{Editorial note.} All delivered mathematical text is the author's own, copied from the pinned Stacks Project chapter file sources/stacks-spaces-simplicial/pinned-original-source.tex; no statement or proof step was rewritten, reordered or removed.\par\medskip
\noindent\textit{Editorial note.} The counted proof is the author's complete argument: conditions (2) and (3) are delegated to two dedicated lemmas of the same chapter (lemma-descent-qcoh-for-fppf-hypercovering and lemma-cohomological-descent-qcoh-for-fppf-hypercovering) whose statements are verbatim the required criterion conditions - the author's modular structure, not an omission; the quasi-coherent characterization used is delivered as card L46152 of this batch.\par\medskip
\noindent\textit{Editorial note.} The label lemma-coh-descent-qcoh-unbounded-for-fppf-hypercovering names the counted statement, and the proof environment immediately following it in the source is delivered verbatim as its proof.\par\medskip
\noindent\textit{Editorial note.} The counted statement is carried by the author's own lemma environment in the chapter source.\par\medskip
\noindent\textit{Allowed standard prerequisites (editorial).}\begin{itemize}
\item P-coverings are flat, finitely presented and surjective. Hypercoverings chapter, Remark \textbackslash{}label\{hypercovering-remark-P-covering\}. Lines 7679-7682: the maps a\_n : U\_n -> X and d\_i\textasciicircum{}n : U\_n -> U\_\{n-1\} are flat, locally of finite presentation, surjective.
\item Quasi-coherent = cartesian with quasi-coherent components on the total site. Same chapter, \textbackslash{}label\{lemma-quasi-coherent-sheaf\} (card L46152 of this batch). Lines 7684-7687: the characterization of QCoh(O\_U) inside Mod(O\_U).
\item Cartesian modules form a weak Serre subcategory. Same chapter, \textbackslash{}label\{lemma-Serre-subcat-cartesian-modules\}. Lines 7688-7691: with flatness of the d\_i\textasciicircum{}n shown above.
\item QCoh on an algebraic space is a weak Serre subcategory. Properties of Spaces chapter, \textbackslash{}label\{spaces-properties-lemma-properties-quasi-coherent\}. Lines 7692-7695: the component statement combined into QCoh(O\_U) weak Serre.
\item Cohomology on Sites, the unbounded equivalence criterion. Cohomology on Sites chapter, \textbackslash{}label\{sites-cohomology-lemma-equivalence-unbounded-one\}. Lines 7700-7704: conditions (1)-(5) checked one by one.
\item Flatness of the augmentation. Same chapter, \textbackslash{}label\{lemma-flat-augmentation-modules\}. Line 7708: condition (1) from flatness of the a\_n.
\item Descent of quasi-coherent modules for fppf hypercoverings. Same chapter, \textbackslash{}label\{lemma-descent-qcoh-for-fppf-hypercovering\}. Lines 7711-7713: condition (2).
\item Cohomological descent for quasi-coherent modules. Same chapter, \textbackslash{}label\{lemma-cohomological-descent-qcoh-for-fppf-hypercovering\}. Lines 7716-7718: condition (3).
\item Cohomology on the total site computes componentwise. Same chapter, \textbackslash{}label\{lemma-sanity-check-modules\}. Lines 7745-7747: H\textasciicircum{}p(V/U\_n, F) = H\textasciicircum{}p(V, F\_n) for the vanishing in condition (4).
\item Quasi-coherent cohomology vanishes on affines. Cohomology of Schemes chapter, \textbackslash{}label\{coherent-lemma-quasi-coherent-affine-cohomology-zero\}; Cohomology of Spaces, Section \textbackslash{}ref\{spaces-cohomology-section-higher-direct-image\}. Lines 7749-7754: H\textasciicircum{}p(V, F\_n) = 0 for p > 0 with V affine.
\end{itemize}\medskip
\ifdefined\hypertarget\hypertarget{sourcepointer-1}{}\fi
\noindent\textit{Source reference.} Original lemma 01XB: sources/stacks-chow/tags.txt, source lines 2512--2512. Prerequisite pointer invoked at proof line(s) 7758; it names a standard fact of a sibling chapter that the counted argument uses, and every step of the counted proof built on it is retained verbatim in the delivered body.\par\medskip
\ifdefined\hypertarget\hypertarget{sourcepointer-2}{}\fi
\noindent\textit{Source reference.} Original remark 0DER: sources/stacks-chow/tags.txt, source lines 16811--16811. Prerequisite pointer invoked at proof line(s) 7693; it names a standard fact of a sibling chapter that the counted argument uses, and every step of the counted proof built on it is retained verbatim in the delivered body.\par\medskip
\ifdefined\hypertarget\hypertarget{sourcepointer-3}{}\fi
\noindent\textit{Source reference.} Original reference 0D7J: sources/stacks-spaces-simplicial/pinned-original-source.tex, source lines 2669--2669. Prerequisite pointer invoked at proof line(s) 7699; it names a standard fact of the same chapter that the counted argument uses, and every step of the counted proof built on it is retained verbatim in the delivered body.\par\medskip
\ifdefined\hypertarget\hypertarget{sourcepointer-4}{}\fi
\noindent\textit{Source reference.} Original reference 0DHE: sources/stacks-spaces-simplicial/pinned-original-source.tex, source lines 7628--7628. Prerequisite pointer invoked at proof line(s) 7726; it names a standard fact of the same chapter that the counted argument uses, and every step of the counted proof built on it is retained verbatim in the delivered body.\par\medskip
\ifdefined\hypertarget\hypertarget{sourcepointer-5}{}\fi
\noindent\textit{Source reference.} Original reference 0DHD: sources/stacks-spaces-simplicial/pinned-original-source.tex, source lines 7544--7544. Prerequisite pointer invoked at proof line(s) 7722; it names a standard fact of the same chapter that the counted argument uses, and every step of the counted proof built on it is retained verbatim in the delivered body.\par\medskip
\ifdefined\hypertarget\hypertarget{sourcepointer-6}{}\fi
\noindent\textit{Source reference.} Original reference 0DH3: sources/stacks-spaces-simplicial/pinned-original-source.tex, source lines 2286--2286. Prerequisite pointer invoked at proof line(s) 7718; it names a standard fact of the same chapter that the counted argument uses, and every step of the counted proof built on it is retained verbatim in the delivered body.\par\medskip
\ifdefined\hypertarget\hypertarget{sourcepointer-7}{}\fi
\noindent\textit{Source reference.} Original reference 0D7M: sources/stacks-spaces-simplicial/pinned-original-source.tex, source lines 2817--2817. Prerequisite pointer invoked at proof line(s) 7698; it names a standard fact of the same chapter that the counted argument uses, and every step of the counted proof built on it is retained verbatim in the delivered body.\par\medskip
\ifdefined\hypertarget\hypertarget{sourcepointer-8}{}\fi
\noindent\textit{Source reference.} Original reference 0DH2: sources/stacks-spaces-simplicial/pinned-original-source.tex, source lines 2208--2208. Prerequisite pointer invoked at proof line(s) 7751; it names a standard fact of the same chapter that the counted argument uses, and every step of the counted proof built on it is retained verbatim in the delivered body.\par\medskip
\ifdefined\hypertarget\hypertarget{sourcepointer-9}{}\fi
\noindent\textit{Source reference.} Original section 0DE7: sources/stacks-spaces-simplicial/pinned-original-source.tex, source lines 6985--6985. Prerequisite pointer invoked at proof line(s) 7687; it names a standard fact of the same chapter that the counted argument uses, and every step of the counted proof built on it is retained verbatim in the delivered body.\par\medskip
\ifdefined\hypertarget\hypertarget{sourcepointer-10}{}\fi
\noindent\textit{Source reference.} Original lemma 0D7V: sources/stacks-chow/tags.txt, source lines 16570--16570. Prerequisite pointer invoked at proof line(s) 7714; it names a standard fact of a sibling chapter that the counted argument uses, and every step of the counted proof built on it is retained verbatim in the delivered body.\par\medskip
\ifdefined\hypertarget\hypertarget{sourcepointer-11}{}\fi
\noindent\textit{Source reference.} Original reference 0D6R: sources/stacks-chow/tags.txt, source lines 16531--16531. Prerequisite pointer invoked at proof line(s) 7733; it names a standard fact of a sibling chapter that the counted argument uses, and every step of the counted proof built on it is retained verbatim in the delivered body.\par\medskip
\ifdefined\hypertarget\hypertarget{sourcepointer-12}{}\fi
\noindent\textit{Source reference.} Original section 071Y: sources/stacks-chow/tags.txt, source lines 9007--9007. Prerequisite pointer invoked at proof line(s) 7756; it names a standard fact of a sibling chapter that the counted argument uses, and every step of the counted proof built on it is retained verbatim in the delivered body.\par\medskip
\ifdefined\hypertarget\hypertarget{sourcepointer-13}{}\fi
\noindent\textit{Source reference.} Original lemma 03M1: sources/stacks-chow/tags.txt, source lines 4733--4733. Prerequisite pointer invoked at proof line(s) 7706; it names a standard fact of a sibling chapter that the counted argument uses, and every step of the counted proof built on it is retained verbatim in the delivered body.\par\medskip
\begin{lemma}
\label{lemma-coh-descent-qcoh-unbounded-for-fppf-hypercovering}
Let $S$ be a scheme. Let $X$ be an algebraic space over $S$.
Let $U$ be a simplicial algebraic space over $S$. Let $a : U \to X$
be an augmentation. Assume $a : U \to X$ is an fppf hypercovering of $X$.
Then $\QCoh(\mathcal{O}_U)$ is a weak Serre subcategory of
$\textit{Mod}(\mathcal{O}_U)$ and
$$
a^* : D_\QCoh(\mathcal{O}_X) \longrightarrow D_\QCoh(\mathcal{O}_U)
$$
is an equivalence of categories with quasi-inverse given by
$Ra_*$. Here $a : \Sh(U_\etale) \to \Sh(X_\etale)$
is as in Section \ref{section-simplicial-algebraic-spaces}.
\end{lemma}

\begin{proof}
First observe that the maps $a_n : U_n \to X$ and $d^n_i : U_n \to U_{n - 1}$
are flat, locally of finite presentation, and surjective by
Hypercoverings, Remark \ref{hypercovering-remark-P-covering}.

\medskip\noindent
Recall that an $\mathcal{O}_U$-module $\mathcal{F}$ is quasi-coherent if and
only if it is cartesian and $\mathcal{F}_n$ is quasi-coherent for all $n$.
See Lemma \ref{lemma-quasi-coherent-sheaf}.
By Lemma \ref{lemma-Serre-subcat-cartesian-modules}
(and flatness of the maps $d^n_i : U_n \to U_{n - 1}$ shown above)
the cartesian modules for a weak Serre subcategory of
$\textit{Mod}(\mathcal{O}_U)$. On the other hand
$\QCoh(\mathcal{O}_{U_n}) \subset \textit{Mod}(\mathcal{O}_{U_n})$
is a weak Serre subcategory for each $n$
(Properties of Spaces, Lemma
\ref{spaces-properties-lemma-properties-quasi-coherent}).
Combined we see that
$\QCoh(\mathcal{O}_U) \subset \textit{Mod}(\mathcal{O}_U)$
is a weak Serre subcategory.

\medskip\noindent
To finish the proof we check the conditions (1) -- (5) of
Cohomology on Sites, Lemma
\ref{sites-cohomology-lemma-equivalence-unbounded-one} one by one.

\medskip\noindent
Ad (1). This holds since $a_n$ flat (seen above) implies $a$ is flat
by Lemma \ref{lemma-flat-augmentation-modules}.

\medskip\noindent
Ad (2). This is the content of
Lemma \ref{lemma-descent-qcoh-for-fppf-hypercovering}.

\medskip\noindent
Ad (3). This is the content of
Lemma \ref{lemma-cohomological-descent-qcoh-for-fppf-hypercovering}.

\medskip\noindent
Ad (4). Recall that we can use either the site $U_\etale$ or
$U_{spaces, \etale}$ to define the small \'etale topos
$\Sh(U_\etale)$, see Section \ref{section-simplicial-algebraic-spaces}.
The assumption of
Cohomology on Sites, Situation \ref{sites-cohomology-situation-olsson-laszlo}
holds for the triple
$(U_{spaces, \etale}, \mathcal{O}_U, \QCoh(\mathcal{O}_U))$
and by the same reasoning for the triple
$(U_\etale, \mathcal{O}_U, \QCoh(\mathcal{O}_U))$.
Namely, take
$$
\mathcal{B} \subset \Ob(U_\etale) \subset \Ob(U_{spaces, \etale})
$$
to be the set of affine objects. For $V/U_n \in \mathcal{B}$
take $d_{V/U_n} = 0$ and take $\text{Cov}_{V/U_n}$ to be the set of
\'etale coverings $\{V_i \to V\}$ with $V_i$ affine.
Then we get the desired vanishing because for
$\mathcal{F} \in \QCoh(\mathcal{O}_U)$
and any $V/U_n \in \mathcal{B}$ we have
$$
H^p(V/U_n, \mathcal{F}) = H^p(V, \mathcal{F}_n)
$$
by Lemma \ref{lemma-sanity-check-modules}. Here on the
right hand side we have the cohomology of the quasi-coherent
sheaf $\mathcal{F}_n$ on $U_n$ over the affine object $V$
of $U_{n, \etale}$. This vanishes for $p > 0$ by the discussion in
Cohomology of Spaces, Section
\ref{spaces-cohomology-section-higher-direct-image} and
Cohomology of Schemes, Lemma
\ref{coherent-lemma-quasi-coherent-affine-cohomology-zero}.

\medskip\noindent
Ad (5). Follows by taking $\mathcal{B} \subset \Ob(X_{spaces, \etale})$
the set of affine objects and the references given above.
\end{proof}

\end{document}
